\interlude{牛奶之尸}

《魔法少女小圆》中最突出的主题之一，是衰败。
熵显然是衰败的一种形式，因此，魔法少女／魔女被塑造成了对抗衰败的武器。
然而，作品中还有其他形式的衰败在不断发生：
城市随着剧情推进日渐破败，从第一话中明亮洁净的空间，一步步变成第十二话中摇摇欲坠的废墟。
最显著的，则是魔法少女自身精神状态的衰败。
这一点在沙耶香的沉沦中表现得最为明显，但作品也给出了大量暗示，表明其他魔法少女同样深受抑郁折磨：
麻美与小圆交心时，画面中出现了一座看似由百忧解\footnote{%
    译注：
    作者指动画第三集 16:35 至 16:45 左右时，出现的大量蓝色胶囊从画面上方倾泻而下的背景意象。
    百忧解是一种抗抑郁药物，学名氟西汀。

}构成的喷泉；
杏子则总是靠不断进食来排遣情绪。
魔女系统的全部目的，就是让魔法少女在情感上不断衰败，直至变成魔女。
从某种意义上说，丘比的系统所做的一切，不过是把熵从宇宙的物质衰败转移成少女们的情感衰败。

衰败始终萦绕全篇，与动画的佛教思想根基紧密相连。
佛教核心教义“四圣谛”中的第一谛指出，“苦”（dukkha）无可避免；这个词大致可以理解为痛苦。
苦有三种：
疾病、衰老和死亡带来的普通而显而易见的苦；
试图执守受时间支配、因而不断变化的事物时产生的焦虑之苦；
以及一切物质性事物因其转瞬即逝而内含的根本之苦\cite{ref108}。

最后一种苦大致可以对应于熵，也就是一切物质性事物最终必然走向衰竭的原则\cite{ref109}。
这种无可避免的衰败看似只会带来绝望，但它并非没有解决办法。
佛教给出的主要方案是不执著——
脱离此世，也就是脱离必然通向绝望的业力轮回\cite{ref110}。
在动画结尾，小圆以观音菩萨的身份，为魔法少女们打开的正是这扇门。
但这是唯一的解决办法吗？
在这个无常的世界中，难道就没有获得幸福的可能吗？

西方文化最初给出的答案同样是“没有”。
基督教的解决办法，是离开这个世界、升入天堂；
此外，它还认为，在未来的某个时刻，上帝会毁灭这个充满苦难的世界，代之以一个更好的世界。
然而，中世纪与文艺复兴时期出现了另一个概念，为人们提供了摆脱衰败与绝望的不同道路：
腐化。

“腐化”原是一个炼金术概念，也是发酵的另一种称呼，后来则逐渐被用来描述死亡与腐烂如何孕育生命\cite{ref111}。
设想一枚正在腐烂的果实，它令人的感官不适：
漆黑、丑陋，而且散发着恶臭，但与此同时，霉菌与蛆虫等新生命却在其上恣意迸发。
这些生命又会成为“更高等”生命形式的养料，到头来，就连最高等的生物也必须依靠腐烂才能生存。

这不仅是生物学意义上的生命循环，也是炼金术士最深刻的精神教诲之一：
死亡孕育生命。
腐烂与创造本为一体。
衰败即是进化。

或者换一种说法：
花朵会在墓园中开放。
彼岸花便是其中之一；
《叛逆》中，焰的魔女形态头戴的花冠正由这种花构成。
彼岸花多呈鲜艳的红色，与大多数花不同，它开花之前，叶片便已凋落。
它常与那些因命运和死亡而分离的挚爱联系在一起，在日本也经常被种植于墓园\cite{ref112}。
这与焰的痛苦显然彼此呼应：
她被迫同挚爱的小圆分离。

然而，种下这些花，意味着人仍未忘记逝去的挚爱；
即便物质存在已经腐朽消散，但爱依然能够延续。
事实上，正是这份源自一个已毁灭宇宙的爱，把小圆带回了焰的虚幻世界。
小圆还带来了另外两个存在：
夏洛特与沙耶香。
两者固然都是出于对小圆的责任与忠诚而归来，却也在后来道出了各自另一层动机。

考虑到夏洛特在整部《叛逆》中一直痴迷于奶酪，她回来是为了奶酪，这并不意外。
奶酪是腐化的绝佳象征：
它美味而富有营养，本质上却又是正在腐烂的牛奶。
夏洛特并不是唯一抱着这种动机归来的角色；
小圆与她的两名随从之所以归来，都是为了某种从衰败中诞生的宝物。
对小圆而言，这种事物是她与焰之间历经多条时间线才逐渐发展起来的关系。
在此过程中，小圆从一个明朗快活的魔法少女逐渐衰败，成为本篇所处时间线中那个近乎全然被动的人物，焰则从一个怯懦的少女逐渐变得彻底封闭。
另一方面，沙耶香是为了她与杏子的关系而来，这段关系源于沙耶香的精神状态迅速衰败时，杏子曾试图走近她。

换言之，腐化的产物可能具有价值。
衰败并非纯然的恶。

那么，我们又该如何看待这一事实：
小圆那完美无瑕、没有衰败的涅槃之中，必然不存在奶酪？
一方面，那里根本无法制造奶酪；
另一方面，如果那里有奶酪，夏洛特也就没有理由离开。
在一个没有衰败、没有腐化的世界中，一切由腐烂创造的美与生命都不可能存在。
小圆的世界中既没有奶酪，也没有由烈火淬炼而成的友谊，因此，它不能算是解决衰败问题的充分方案。

焰给出的答案是否更好，只能留待时间及《叛逆》那看来不可避免的续作来回答了。
